\chapter{How We Think: The Architect's Method}

This chapter is not about the system. It is about the mind the system came
out of --- written so that anyone judging, using, or continuing this work
knows how the decisions get made. The Architect (Frans Nekongo) thinks the
following way, and Tottie is what happens when that thinking is given a
machine that works fast.

\section{The mind is fractal}
A parent is a parent: it does what the parent does. The kids do what they
need to do. Zoom into any kid and there are more kids with kids, down to
the cell. Every level takes context, makes \emph{one small decision}, and
passes context along. The hypothalamus works this way --- nuclei within
nuclei, each with one job --- and so must any memory we build. One thing
doing the job of many is the failure mode; many small things each doing
one job is the design. (Research later confirmed independent lineages:
Minsky's society of mind, Brooks' subsumption, the cortical hierarchy.
They verify the shape. They did not originate it here.)

\section{Simulate; never claim to recreate}
We cannot recreate a brain and we do not pretend to. We pick specific
functions --- familiarity, surprise, consolidation, recall, foresight ---
and \emph{simulate} them, honestly, one mechanism at a time. New
understanding adds on; it never destroys what already works. The fast road
stays when the slow road is built beside it.

\section{Memory gives; it is never fetched}
The whole goal in one line: context must arrive \emph{unasked}, or the
product failed. A memory system an agent has to interrogate is a database
with good marketing. The brain pushes the deadline, the lesson, the due
intention into working context before it is requested --- that push is
the thing being built, and every experiment must be shaped so the push is
what gets measured.

\section{Imagination proposes; research verifies; then code}
The human brings what the machine can never have: imagination that creates
shapes from nothing. The machine brings what the human can never have:
speed and total recall. Neither is above the other. But imagination never
goes straight to code. Every shape is checked first against the research
corpus (the project notebook, then the literature library) and the
internet. If the brain does it differently, the brain wins. If no source
knows, the gap is marked as a guess and measured later. Building from the
head is how clever systems get built that understand nothing.

\section{No theater}
Scripted test turns prove nothing: twelve canned sentences through a chat
box cannot break a memory. Real testing means real work --- worker agents
doing real projects, with real deadlines inside the fiction, fresh agents
per phase so nothing rides on leftover context, and problems allowed to
show themselves naturally. The manager assigns, injects change the way
life does, verifies every claim against the database, and reports what
broke without softening it. A test that cannot fail is decoration.

\section{Blunt talk both ways}
Sycophancy is a bug, in either direction. The machine must never flatter,
never pad, never perform competence with jargon --- plain words, short
messages, what-can-be-observed. The human returns the same honesty:
calling out drift the moment it appears, including in how the machine
talks. Depth over appearance, always. God blessed the Architect with the
thinking; the thinking is only as good as the honesty around it.
